\chapter{Multiblocks}
\label{ch:multiblocks}

Every structure here is built by placing its blocks in the right shape and then striking it with an
Engineer's Hammer. The in-game manual renders each one as a rotating diagram; \autoref{ch:ingame}
carries the text, and the exact block lists are in \autoref{ch:blocks}.

\begin{ietip}
Build multiblocks with the hammer already in your hand and the manual page open. Almost every failed
assembly is one block in the wrong place, and the diagram will show you which.
\end{ietip}

\section{Processing}

\begin{description}
\item[Coke Oven] $3\times3\times3$. Coal to coal coke and creosote oil. Needs no power.
\item[Blast Furnace] $3\times3\times3$. Iron to steel, with slag. Needs no power.
\item[Advanced Blast Furnace] Faster, and accepts a Preheater for more speed still.
\item[Alloy Smelter] Combines two metals. Needs no power.
\item[Crusher] Doubles ore, and grinds a great many other things. The first machine most players
      build that visibly pays for itself.
\item[Arc Furnace] Recycles, alloys and processes ores further than the Crusher can. Consumes
      graphite electrodes, which wear out.
\item[Metal Press] Ingots to plates, rods and wire with a mould. Cheap to run and enormously
      cheaper than hammering.
\item[Squeezer] Presses seeds to plant oil, and other things besides.
\item[Fermenter] Ferments crops to ethanol.
\item[Refinery] Plant oil and ethanol to biodiesel.
\item[Mixer] Combines fluids with items --- concrete, most importantly.
\item[Bottling Machine] Fills containers with a fluid.
\item[Auto Workbench] Automates blueprint crafting.
\item[Assembler] Automates ordinary crafting from up to three recipes at once, pulling from a
      shared inventory.
\end{description}

\section{Mining}

\begin{description}
\item[Excavator] The largest structure in the mod. Digs a mineral vein located by Sample Drill,
      producing ore in quantity until the vein depletes. Consumes \textbf{4\,096 IF/t} by default.
\item[Bucket Wheel] The visible half of the Excavator, and the reason it is worth building where
      people can see it.
\end{description}

\section{Storage and power}

\begin{description}
\item[Sheetmetal Tank] A $5\times5\times5$ vessel holding 512\,000\,mB.
\item[Silo] The same idea for solid items.
\item[Diesel Generator] \autoref{ch:generators}.
\item[Lightning Rod] \autoref{ch:generators}.
\end{description}

\begin{iefork}
The fork adds several: the Grid Management Console and Substation (\autoref{ch:grid}), the Fluid
Control Console (\autoref{ch:fluidnet}), and the whole petroleum chain --- Derrick, Cracking Unit,
Loading Gantry, the buried tanks and the free-form Storage Tank (\autoref{ch:petroleum}).
\end{iefork}

\begin{iefork}
City mode stops idle multiblocks re-scanning the recipe list every tick and widens the scan interval
for the rest. A room full of idle machines is close to free.
\end{iefork}
